% BHU PYQ -- Manually transcribed & error-corrected
% Paper: GLB-602: Stratigraphy (Special Examination)
% Exam: B.Sc. (Hons.) Semester IV (Special) Examination, 2023-24

\documentclass[11pt,a4paper]{article}
\usepackage[a4paper,top=2.5cm,bottom=2.5cm,left=2.8cm,right=2.8cm,headheight=14pt]{geometry}
\usepackage{amsmath,amssymb}
\usepackage[T1]{fontenc}
\usepackage[utf8]{inputenc}
\usepackage{lmodern,microtype,fancyhdr,enumitem,hyperref,lastpage,parskip}

\hypersetup{colorlinks=true,urlcolor=black,linkcolor=black,
  pdftitle={GLB-602 Stratigraphy (Special) -- BHU B.Sc. Sem IV (Special) 2023-24}}
\pagestyle{fancy}\fancyhf{}
\renewcommand{\headrulewidth}{0.4pt}\renewcommand{\footrulewidth}{0.4pt}
\fancyhead[L]{\small   \textit{Banaras Hindu University}}
\fancyhead[R]{\small   \textit{GLB-602: Stratigraphy (Special)}}
\fancyfoot[C]{\small Page  \thepage\ of \pageref{LastPage}}


\newlist{parts}{enumerate}{2}
\setlist[parts,1]{label=(\alph*),leftmargin=*,itemsep=5pt,topsep=3pt}

\newcommand{\pts}[1]{\hfill   \textit{[#1]}}


\begin{document}

\begin{center}
  {\large\bfseries BANARAS HINDU UNIVERSITY}\\[2pt]
  {\normalsize B.Sc. (Hons.) --- Semester IV (Special) Examination, 2023-24}\\[6pt]
  \rule{0.8\linewidth}{0.5pt}\\[6pt]
  {\large\bfseries Paper: GLB-602 --- Stratigraphy (Special Examination)}\\[4pt]
  {\normalsize Time: 3 Hours \hfill Maximum Marks: 70}\\[2pt]
  {\small Ref. No.: 077021}
\end{center}
\rule{\linewidth}{0.4pt}
\smallskip



\noindent   \textit{Instructions: Answer   \textbf{five} questions in all, selecting at least   \textbf{two} each from Sections A and B, and Section-C which is compulsory. All questions carry equal marks.}
\bigskip

\bigskip


\noindent {\large\bfseries SECTION --- A}
\rule{\linewidth}{0.4pt}\smallskip




\noindent   \textbf{Question 1.} Write a brief note on the `Principles of Stratigraphy' and discuss the various criteria for stratigraphic classification. \pts{14}

\medskip


\noindent   \textbf{Question 2.} Write short notes on   \textbf{any two} of the following: \pts{$2  \times 7 = 14$}

\begin{parts}
  \item Accommodation space in sequence stratigraphy
  \item Ariyalur Group
  \item Stratigraphic correlation
\end{parts}

\medskip


\noindent   \textbf{Question 3.} Give an account of the lithostratigraphic framework and fossil content of the Mesozoic succession of the Kachchh Basin, Western India. \pts{14}

\medskip


\noindent   \textbf{Question 4.} Write in brief on   \textbf{any two} of the following: \pts{$2  \times 7 = 14$}

\begin{parts}
  \item Uttatur Group
  \item Palaeozoic of Kashmir
  \item Upper Triassic succession of Spiti
\end{parts}

\bigskip


\noindent {\large\bfseries SECTION --- B}
\rule{\linewidth}{0.4pt}\smallskip




\noindent   \textbf{Question 5.} Discuss the stratigraphy of the marine Palaeogene succession of India. \pts{14}

\medskip


\noindent   \textbf{Question 6.} Briefly discuss the origin, classification, and age of the Deccan Traps. \pts{14}

\medskip


\noindent   \textbf{Question 7.} Discuss briefly   \textbf{any two} of the following: \pts{$2  \times 7 = 14$}

\begin{parts}
  \item Upper Gondwana
  \item Neogene -- Quaternary stratigraphic succession of the Siwalik Supergroup
  \item Lower Triassic of Spiti
\end{parts}

\bigskip


\noindent {\large\bfseries SECTION --- C}
\rule{\linewidth}{0.4pt}\smallskip




\noindent   \textbf{Question 8.} Define the following in a maximum of five sentences each: \pts{$7  \times 2 = 14$}

\begin{parts}
  \item Type section
  \item Muth Formation
  \item Erathem
  \item Anaipadi Formation
  \item Depositional sequence
  \item Intertrappean beds
  \item Fenestella Shales
\end{parts}

\vfill\rule{\linewidth}{0.4pt}

\begin{center}\small
\href{https://bhu-exam.vercel.app/}{Click here to visit our website}\\[4pt]
\href{https://me-aryan.vercel.app/}{Click here to visit our contributor}\\[4pt]
\href{https://quashan-me.vercel.app/}{Click here to visit our contributor}
\end{center}
\end{document}
